\chapter{SimulationCode}\label{ch:simulation-code}
\chapterstatus{written}{2026-11-20}

SimulationCode is the complete particle--mesh and COLA simulation program of \fml{}: one parameter file sets up the
initial conditions, the evolution and the analysis during the run. It uses the slab decomposition, so it is not meant
for very high resolution, but is well suited to fast approximate simulations and large suites of runs, to testing
new models and methods, and to analysing snapshots. This chapter describes the workflow, the matter content
(species, massive neutrinos, the $N$-body gauge), the outputs and on-the-fly analysis, and the validation. The physics
of the steps is in \cref{ch:growth-transfer,ch:random-fields,ch:lpt,ch:nbody-pm,ch:short-range}, the cosmologies and
gravity models in \cref{ch:mg-simulations}.

\begin{modulebox}{SimulationCode}
Program \code{nbody} (\file{src/Main.cpp}); the simulation \file{src/Simulation.h}, \file{SimulationInit.h},
\file{SimulationStep.h}, \file{RunSimulation*.h/.cpp} (instantiations for the particle types), \file{KickDrift.h},
\file{COLA.h}; models \file{Cosmology.h}, \file{GravityModel*.h}; matter \file{Species.h}, \file{Neutrinos.h},
\file{NeutrinoParticles.h}, \file{NBodyGauge.h}, \file{SimulationIDE.h}; IC \file{InitialDensityField.h},
\file{TransferInput.h/.cpp}, \file{ExternalIC.h}; analysis \file{AnalyzeOutput.h}, \file{HaloFinding.h},
\file{SnapshotOutput.h}, \file{Lightcone/}; parameters \file{SimulationParameters.h}, \file{ParameterSchema.h}.\par
Documentation: \file{SimulationCode/README.md}, \file{PARAMETERS.md}, \file{MODELS.md}, \file{INITIAL_CONDITIONS.md},
\file{EVOLUTION.md}, \file{SPECIES.md}, \file{ANALYSIS.md}, \file{DEVELOPING.md}, \file{tests/README.md}; an
annotated parameter file \file{parameterfile.lua}.
\end{modulebox}

\section{Workflow}\label{sec:sim-workflow}

A run proceeds as follows.
\begin{enumerate}
\item \textbf{Parameters} are read from a Lua file and checked against a schema (\file{ParameterSchema.h}): types,
      allowed values, dependencies (``used when''), incompatible combinations (an error) and unknown names (a warning
      suggesting the intended name). \code{./nbody --parameters PARAMETERS.md} writes the reference of all parameters
      (also \cref{app:simcode-parameters}). Parameters not in the file take their defaults, unless
      \param{all_parameters_must_be_in_file} is set: then every parameter that is used must be given, so the file
      documents the run completely.
\item \textbf{Cosmology and gravity model}: the background from the registry of \cref{ch:background}
      (\param{cosmology_model}; besides the densities every cosmology has the curvature \param{cosmology_OmegaK}, the CMB
      temperature \param{cosmology_TCMB_kelvin} and \param{cosmology_Neffective}) and the gravity model (\param{gravity_model}), which provides
      $G_{\rm eff}(k,a)/G$, the growth factors (\cref{ch:growth-transfer}) and the force.
\item \textbf{Transfer input} (\param{ic_type_of_input}: a $P(k)$ or $T(k)$ file, \code{camb}, \code{class},
      \code{internal}, \code{tables}; \cref{sec:transfer}), with optional $\sigma_8$ normalisation
      (\param{ic_sigma8_normalization}: \param{ic_sigma8} at \param{ic_sigma8_redshift}). CLASS output may be in either gauge
      (\param{ic_input_gauge}; the Newtonian gauge is converted to the CDM-synchronous gauge of the initial conditions).
      For modified gravity \param{ic_use_gravity_model_GR} treats the input as a $\Lambda$CDM spectrum and scales it to
      the initial redshift with the GR growth, so the initial conditions are those of the matching $\Lambda$CDM run.
\item \textbf{Initial conditions}: a Gaussian or non-Gaussian field on an \param{ic_nmesh}$^3$ grid
      (\cref{ch:random-fields}), LPT displacements to order \param{ic_LPT_order} (\cref{ch:lpt}) for each species; or
      particles or phases from a GADGET snapshot (\param{ic_reconstruct_gadgetfilepath}), whose density field is
      reconstructed with \param{ic_reconstruct_assigment_method} and \param{ic_reconstruct_interlacing} and smoothed with
      \param{ic_reconstruct_smoothing_filter} on the scale \param{ic_reconstruct_dimless_smoothing_scale} (in units of the
      box).
\item \textbf{Evolution}: fixed steps between the output redshifts (\param{timestep_nsteps},
      \param{timestep_method}; uniform in $a^p$ with \param{timestep_spacing_power} for the \code{powerlaw} spacing,
      \cref{sec:nbody-timestepping}; kick--drift--kick, \param{timestep_algorithm}, the only ordering implemented); in each
      step the density (and its $P(k)$, almost free since it is already in Fourier
      space), the force of the gravity model, the kick and the drift, with COLA if enabled (with
      \param{simulation_enforce_LPT_trajectories} the PM force is skipped and the particles follow their LPT trajectories,
      a check of the COLA frame); particle migration between slabs; the lightcone.
\item \textbf{Outputs} at \param{output_redshifts}: snapshots (\param{output_particles}; GADGET binary or HDF5,
      \param{output_fileformat}) and the analyses switched on.
\end{enumerate}

\begin{paramfile}[A minimal COLA run (excerpt of parameterfile.lua)]
simulation_name      = "TestSim"
simulation_boxsize   = 512.0           -- Mpc/h
simulation_use_cola  = true
cosmology_model      = "LCDM"
cosmology_OmegaCDM   = 0.25;  cosmology_Omegab = 0.05;  cosmology_h = 0.7
cosmology_As = 2.1e-9;  cosmology_ns = 0.966;  cosmology_kpivot_mpc = 0.05
gravity_model        = "GR"
particle_Npart_1D    = 128
ic_type_of_input     = "internal";  ic_input_redshift = 20.0;  ic_initial_redshift = 20.0
ic_LPT_order         = 2;  ic_random_seed = 1234;  ic_fix_amplitude = true
timestep_method      = "Quinn";  timestep_nsteps = {30}
force_nmesh          = 256
output_redshifts     = {0.0};  output_folder = "output"
pofk = true;  fof = true
\end{paramfile}
\begin{paramfile}[Running]
mpirun -np 4 ./nbody parameterfile.lua
./nbody --write-transfers parameterfile.lua tables.txt   # the solver's transfer tables, once per cosmology
\end{paramfile}

\paragraph*{Particle types.} By default one particle type that can run every configuration is compiled (120 bytes
in 3D). With \code{USE_MINIMAL_PARTICLES} all types are compiled and the smallest one that a run needs is used
(56--104 bytes; 120 with masses): position, velocity, ID, the stored LPT displacement fields for COLA, the Lagrangian
coordinate, the mass and the neutrino weight are fields that a type has or not (compile-time reflection,
\cref{ch:particles}).

\section{Matter content}\label{sec:sim-species}

\paragraph*{CDM and baryons.} By default CDM and baryons are one species with the CDM+baryon transfer function. With
\param{particle_baryons} they are two sets of particles with masses in the ratio $\Om_b/\Om_{\rm cdm}$ that feel the
total CDM+baryon density, with initial conditions from one random field and the transfer-function ratios
(\cref{sec:linear-input}): each species displaced by its own linear field (\code{displacement}), or the relative mode
carried by the particle masses (\code{mass}, recommended, or \code{mass_ini}; \param{particle_baryon_relative_mode}).
Both species start from the same lattice (\param{particle_baryon_lattice} \code{= same}): interleaved lattices give a
large spurious relative mode in a PM code (+8\% at $k=0.4\hMpc$ in linear theory). With the same lattice the
two-species CDM+baryon $P(k)$ equals the one-species one, and the relative mode follows linear theory. In COLA both
species share the frame of the total CDM+baryon field; Tassev's modified operators get the relative velocity wrong
($+30\%$ in $P_{b-c}$ with 20 linear steps), so they are not allowed with two species.

\paragraph*{Neutrino masses.} Three species of temperature $T_\nu=T_{\rm CMB}(N_{\rm eff}/3)^{1/4}(4/11)^{1/3}$:
three equal masses from \param{cosmology_OmegaMNu}, or individual masses \param{cosmology_neutrino_masses_eV}. The
background treats them exactly (\cref{sec:bg-ncdm}).

\paragraph*{Linear massive neutrinos} (\param{force_linear_massive_neutrinos}). The neutrinos are a linearly evolved
field: the force is sourced by
\begin{equation}
\delta_m(\vk,a)=\left[(1-f_\nu)+f_\nu\frac{T_\nu(k,a)}{T_{cb}(k,a)}\right]\delta_{cb}(\vk,a),\qquad f_\nu=\frac{\Om_\nu}{\Om_m},
\end{equation}
with the transfer functions of the provider, and the growth factors are scale dependent (scale-dependent COLA).
With \param{force_neutrino_energy_density} the neutrinos source the Poisson equation with their energy density,
$f_\nu\langle E/m\rangle(a)\,\delta_\nu$ with $\langle E/m\rangle=\langle\sqrt{1+(q/am)^2}\rangle$ over the Fermi--Dirac
distribution, instead of their mass density (for linear neutrinos and for neutrino particles, which are then weighted by
$E/m$); it matters at high redshift for light neutrinos ($\langle E/m\rangle=1.33$ at $z=50$ for 0.033\,eV).

\paragraph*{Starting at high redshift with massive neutrinos.} While they are relativistic the neutrinos are radiation,
as massless neutrinos in $\Lambda$CDM. If they are relativistic at the start ($\langle E/m\rangle>3$; $z\gtrsim200$ for
0.033\,eV), the growth factors start from the radiation--matter growing mode instead of the growth rate of the transfer
tables (which hold the synchronous-gauge super-horizon mode there), the neutrinos do not cluster before the tables
begin (in the growth factors and in the force), and back-scaled initial conditions only need the tables from where the
neutrinos become non-relativistic. Converged runs from $z=1000$ and $z=10^7$ then agree with a start at $z=50$ to
$\approx0.1\%$ at $z=18$. FastPM and BullFrog take their growth from one scale; FastPM with the scale-dependent kick
(\cref{sec:nbody-timestepping}) is exact for every linear mode.

\paragraph*{Neutrino particles: the $\delta f$ method} \cite{Elbers2021} (\param{particle_neutrinos}). The bulk part of
the initial conditions is made like the CDM, $\delta_\nu=\delta_{cb}T_\nu/T_{cb}$ with the neutrino growth rate; each
particle then gets a thermal momentum $q=y\,k_BT_{\nu0}/c$ in a random direction, with $y$ drawn from the relativistic
Fermi--Dirac distribution $y^2/(e^y+1)$ (\cls{FermiDiracMomentumSampler}). The code velocity is the comoving
momentum $q/(mH_0B)$, and with $u=q/(amc)$ the relativistic equations of motion in the Newtonian gauge are
\begin{equation}
\frac{\dd\vx}{\dd t_{\rm drift}}=\frac{\bm v}{\sqrt{1+u^2}},\qquad
\frac{\dd\bm v}{\dd t_{\rm kick}}=-\bm F\,\frac{1+2u^2}{\sqrt{1+u^2}},
\end{equation}
with the same drift and kick factors as the CDM. The particles carry the $\delta f$ weights
\begin{equation}
w=1-\frac{f_0(y)}{f_i},\qquad f_i=f_0(y_{\rm ini})(1+\delta_{\rm ini}),
\end{equation}
($f_i$ is the phase-space density the particle carries, conserved along its path), so the weighted density
$(N_{\rm grid}^3/N_\nu)\sum_iw_iW(\vx-\vx_i)$ is $\delta_\nu$: the particles carry only the perturbation, which removes
most of the shot noise. The force uses $\delta_m=(1-f_\nu)\delta_{cb}+f_\nu\delta_\nu$. Different masses are one set of
particles, each assigned one of the species by a hash of its ID. There are \param{particle_neutrino_Npart_1D}$^3$ of
them (0: as many as the CDM particles); \param{particle_neutrino_deltaf} \code{= false} drops the weights (plain
particles, much noisier; for tests). Validation ($1024\Mpch$, $128^3$ CDM and $128^3$
neutrino particles): $P_\nu$ follows linear theory to 2\% (0.4\,eV) and 7\% (0.1\,eV) relative to the simulation's own
$P_{cb}$, and $P_{cb}$ agrees with the linear-neutrino run to $<0.2\%$. Not yet with COLA, time stepping other than
Quinn, interlacing, the $N$-body gauge or the lightcone.

\paragraph*{The $N$-body gauge} (\param{force_nbody_gauge}). On very large scales a Newtonian simulation misses the
effect of radiation and neutrinos on the growth, and the meaning of the positions is ambiguous. In the $N$-body gauge
\cite{Fidler2015} the particles follow the Newtonian equations exactly to linear order in GR if the force is sourced
by $\rho_{cb}(\delta_{cb}+\delta_{\rm extra})$ (\cref{sec:transfer-gauges}). The source is added to the density at every
force calculation and the initial conditions use the $N$-body gauge transfer functions; with \code{internal} input
the solver provides both (integrating the linear growth with the source reproduces the solver's $\delta_{cb}$ to
$10^{-5}$), and for CAMB input \file{scripts/make_nbody_gauge_transfer.py} makes them with CLASS
(the table $T_{\rm extra}(k,z)$ given as \param{force_nbody_gauge_filename}). For a start at $z=49$
the effect on $P(k,z=0)$ is $\approx2\%$ at $k=10^{-3}\hMpc$ and $\ll0.1\%$ above $0.01\hMpc$.

\paragraph*{Ultralight axions in the initial conditions.} With \param{ic_axion_fraction} the initial conditions come
from the solver's mixed CDM + axion model (\param{ic_axion_log10_m}, \param{ic_axion_switch_mH}) with the particles
all the dark matter (\code{axion_as_cdm}, \cref{sec:transfer}); the evolution is that of CDM.

\section{On-the-fly analysis and output}\label{sec:sim-analysis}

\begin{itemize}
\item \textbf{Power spectra} of the CDM+baryons and the total matter at every step (\code{pofk_every_step}) with linear
      theory and a nonlinear prediction (EuclidEmulator2, with the e-MANTIS boost for $n=1$ $f(R)$ \cite{SaezCasares2023};
      \cref{ch:fitting-emulators}); at the outputs with a chosen assignment, interlacing and shot-noise subtraction
      (\param{pofk_subtract_shotnoise}, also for the spectra of every step; \cref{ch:power-spectra}); per species with
      separate baryons or neutrino particles.
\item \textbf{Redshift-space multipoles} (plane-parallel, averaged over the three axes; \param{pofk_multipole_nmesh},
      \param{pofk_multipole_density_assignment_method}, \param{pofk_multipole_interlacing},
      \param{pofk_multipole_subtract_shotnoise} for the monopole) and \textbf{bispectra}
      (\param{bispectrum_density_assignment_method}, \param{bispectrum_interlacing}, \param{bispectrum_subtract_shotnoise};
      \cref{sec:pk-multipoles,sec:pk-polyspectra}).
\item \textbf{Hybrid bias spectra} (\code{hybridbias}): the 15 cross spectra of the Lagrangian operators $1$,
      $\delta_L$, $\delta_L^2-\langle\delta_L^2\rangle$, $s^2-\langle s^2\rangle$, $\nabla^2\delta_L$ advected with the
      particles \cite{Modi2020,Kokron2021}, with Zel'dovich and one-loop predictions (\cref{ch:power-spectra}), and with
      \param{hybridbias_rsd} their redshift-space multipoles with the same predictions (\cref{sec:pk-hybrid-rsd}).
\item \textbf{Haloes}: friends-of-friends, relaxed FoF \cite{Dai2019}, spherical-overdensity masses around FoF haloes
      (\code{fof_so_definitions}, e.g.\ \code{200m}, \code{200c}, \code{vir}) and mass functions (\cref{ch:halo-finding}).
      With PM forces haloes are smoothed on the scale of a force cell, so SO masses are biased low unless $R_\Delta$ is
      several force cells (at $1\Mpch$ cells the 200m mass function is 0.56 of the Tinker fit), while FoF masses are much
      less sensitive.
\item \textbf{Baryonification} of a copy of the particles at each output (\cref{ch:pgd-baryonification}).
\item \textbf{Images} (\param{images}): PNG pictures (and optionally \code{.npy} values) of a slab of fields projected
      along an axis at every output, \file{image_<field>_z<z>.png}, made with \file{Visualization}
      (\cref{sec:io-visualization}): \code{density}; \code{velocity_divergence} $-\nabla\cdot\bm v/(aH)$, equal to
      $f\delta$ in linear theory (checked to 2--4\% below $3k_f$); \code{baryons_minus_cdm}; \code{neutrinos} (particles,
      or linear theory on top of $\delta_{cb}$); and \code{fifth_force}, $|\bm F_5|/|\bm F_N|$ with the gravity model's
      own screening method or exact solver (whose state for the next step is kept: the $P(k)$ are the same with and
      without images). Settings \param{images_fields}, \param{images_nmesh}, \param{images_npix}, \param{images_axis},
      \param{images_slab_center}, \param{images_slab_thickness}, \param{images_colormap}, \param{images_npy},
      \param{images_density_assignment_method}.
\item \textbf{Snapshots}: GADGET binary or HDF5 (\param{output_fileformat}; CDM type 1, baryons type 2, neutrinos
      type 3 with their weights).
\item \textbf{Lightcones} (\param{lightcone}): particles crossing the past lightcone of an observer at
      \param{plc_pos_observer} (in units of the box) between \param{plc_z_init} and \param{plc_z_finish}, with box
      replicas (\param{plc_use_replicas}; \param{plc_ndim_rep} of the dimensions replicated in both directions, in 3D 0 for
      an octant and 3 for the full sky), written in ascii, GADGET or HDF5 (\param{plc_output_asciifile},
      \param{plc_output_gadgetfile}, \param{plc_output_hdf5file}), in batches if needed (\param{plc_output_in_batches}),
      and HEALPix ``onion'' density maps $\delta(z,\hat{\bm n})$ (\param{plc_make_onion_density_maps}, resolution
      \param{plc_nside}, shells of width \param{plc_da_maps} in $a$; sparse maps can be stored with chunkpix,
      \param{plc_use_chunkpix} with chunks of \param{plc_nside_chunks}); with \code{plc_relativistic_effects} the
      potential $\Phi/c^2$, radial velocity and the observed redshift $1+z_{\rm obs}=(1+z)(1+v_r/c-\Phi/c^2)$ (Doppler and
      gravitational redshift to first order; no ISW, Shapiro delay or lensing deflection).
\item \textbf{Haloes on the lightcone} (\param{plc_halos}): at every lightcone step friends-of-friends on the particles
      near the shell of the step (within \param{plc_boundary_mpch}, in any replica and across the faces of the box), with
      the linking length \param{plc_halos_linking_length}, at least \param{plc_halos_nmin_per_halo} particles and the
      buffer \param{plc_halos_buffer_mpch} from the neighbouring tasks;
      each halo moves with the mean velocity of its particles and is recorded once, where its centre of mass crosses the
      lightcone, with the same crossing code as the particles (\file{halos_lightcone.txt}: position, velocity, redshift,
      $N$, FoF mass, velocity dispersion). The shell restriction (\param{plc_halos_shell_only}, default) gives the same
      catalogue as linking all the particles.
      Against FoF run afterwards on the particle lightcone ($64^3$ in $128\Mpch$, an octant to $z=0.12$) the halo counts
      per radial shell agree to 2\% and 99.9\% of the haloes with $\geq50$ particles are matched, with 1\% lower masses
      (the haloes are found at the start of the step). In a $192^3$, $1\,h^{-1}$Gpc run with an octant to $z=1$ the haloes
      cost 16\,s of 132\,s.
\end{itemize}
The same analysis layer is available for snapshots of any simulation through SnapshotAnalyzer (\cref{ch:snapshot-analyzer}).

\section{Performance}

Every time-consuming step uses MPI (slabs) and OpenMP. With both, the threads of each task need their own cores (e.g.\
\code{OMP_NUM_THREADS=4 mpirun --map-by slot:PE=4 -np 2 ./nbody params.lua}); the code warns when OpenMPI binds a task
to one core. FFTW plans can be measured (\param{simulation_fftw_planner}) and stored as wisdom, which later runs read with
\param{simulation_fftw_wisdom_file} to skip the planning. Each task allocates \param{particle_allocation_factor} (default
1.5) times its mean number of particles, since particles move between tasks; \param{particle_sort_every_nsteps} reorders
the particles in memory by position every so many steps, which speeds up the density assignment and the force
interpolation once they have moved far from their lattice order (the results change only by round-off). At the end of a run the
time spent in density assignment and FFTs, $P(k)$, force, kick and drift is printed; memory is logged when compiled
with memory logging.

\begin{validation}[Validation (\file{SimulationCode/tests})]
\begin{itemize}
\item \emph{MG-PICOLA}: $128^3$ runs with the same random numbers and time stepping as L-PICOLA/MG-PICOLA agree in
      $P(k,z=0)$ to better than $0.08\%$ up to the Nyquist frequency (the test passes at $0.2\%$).
\item \emph{Linear theory and convergence}: the force grid and gradient kernels (\file{check_force_grid.py}), time
      stepping (\file{check_timestepping.py}), COLA (\file{check_cola.py}), early forces, initial conditions,
      non-Gaussianity, the $N$-body gauge against CLASS, lightcone crossing, haloes.
\item \emph{Transfer input}: CLASS in both gauges and CAMB for the same cosmology agree in $P_{cb}$ and $P_{\rm total}$
      to $1.4\times10^{-3}$ at all outputs; runs from the internal solver agree with CAMB-input runs to $3\times10^{-4}$.
\item \emph{Models}: each cosmology and gravity model against the Einstein--Boltzmann solver's linear growth
      (\cref{ch:mg-simulations}); regression cases with bit-identical CMake and Make builds.
\end{itemize}
Not every combination of models and options has been tested; options that do not work together are listed at the end
of \file{PARAMETERS.md} and give an error.
\end{validation}
